% Source: https://artofproblemsolving.com/community/c3h1509288_new_approximation_for_pi
% Stable identity: aops-topic-1509288-q01
\textbf{Bounding $\dfrac{\sqrt{987}}{10}$.} We have
\[
314165^2=98\,699\,647\,225<98\,700\,000\,000=987\cdot10^{8}<98\,700\,275\,556=314166^2 ,
\]
so $314165<\sqrt{987}\cdot10^4<314166$, that is,
\[
3.14165<\dfrac{\sqrt{987}}{10}<3.14166 .
\]
Since also $3.14159<\pi<3.14160$, the number $\dfrac{\sqrt{987}}{10}$ exceeds $\pi$ by less than $3.14166-3.14159=0.00007$.

\textbf{Bounding $\dfrac{22}{7}$.} Since $7\cdot3.14285=21.99995<22$, we have $\dfrac{22}{7}>3.14285$, so $\dfrac{22}{7}$ exceeds $\pi$ by more than $3.14285-3.14160=0.00125$.

As $0.00007<0.00125$, the approximation $\dfrac{\sqrt{987}}{10}$ is closer to $\pi$ than $\dfrac{22}{7}$. $\blacksquare$

(In fact $\dfrac{\sqrt{987}}{10}=3.1416556\ldots$ is off by about $6.3\cdot10^{-5}$, while $\dfrac{22}{7}$ is off by about $1.26\cdot10^{-3}$, twenty times as much.)
